\documentclass[12pt,a4paper]{article}
\usepackage[utf8]{inputenc}
\usepackage[T1]{fontenc}
\usepackage[brazil]{babel}
\usepackage{amsmath,amssymb}
\usepackage{geometry}
\geometry{margin=2.5cm}
\title{Material de Estudo: Din\^amica Molecular III --- An\'alise e Propriedades}
\author{Princ\'ipio de Modelagem Matem\'atica}
\date{Unidade 29}
\begin{document}
\maketitle

\section*{Objetivo}
Analisar trajet\'orias de DM: fun\c{c}\~ao de distribui\c{c}\~ao radial (RDF), deslocamento quadr\'atico m\'edio (MSD) e coeficiente de difus\~ao.

\section*{Fun\c{c}\~ao de Distribui\c{c}\~ao Radial (RDF)}

\[
g(r) = \frac{\langle n(r) \rangle}{4\pi r^2 \Delta r \cdot \rho},
\]
onde $n(r)$ \'e o n\'umero m\'edio de part\'iculas na casca $[r, r+\Delta r]$ e $\rho = N/V$.
\begin{itemize}
  \item $g(r)=1$ para $r$ grande: gás ideal.
  \item Pico em $r_\text{eq}$: vizinhos mais pr\'oximos.
  \item $g(0) = 0$: part\'iculas n\~ao se sobrep\~oem.
  \item Padr\~ao de picos: s\'olido > l\'iquido > g\'as.
\end{itemize}

\section*{Deslocamento Quadr\'atico M\'edio (MSD)}

\[
\text{MSD}(t) = \frac{1}{N}\sum_{i=1}^N \langle|\vec{r}_i(t) - \vec{r}_i(0)|^2\rangle.
\]

\textbf{Rela\c{c}\~ao de Einstein:}
\[
\text{MSD}(t) = 6Dt \quad \text{(3D, difus\~ao normal)}.
\]
Coeficiente de difus\~ao: $D = \lim_{t\to\infty} \frac{\text{MSD}(t)}{6t}$.

\subsection*{Regimes de difus\~ao}
\begin{itemize}
  \item \textbf{Bal\'istico:} $t$ curto, MSD $\sim t^2$ (movimento livre).
  \item \textbf{Difusivo:} $t$ longo, MSD $\sim t$ (difus\~ao normal).
  \item \textbf{Subdifus\~ao:} MSD $\sim t^\alpha$, $\alpha < 1$ (ambientes confinados).
\end{itemize}

\section*{Rela\c{c}\~ao de Green-Kubo}

$D$ tamb\'em pode ser calculado pela fun\c{c}\~ao de autocorrela\c{c}\~ao de velocidades (VACF):
\[
D = \frac{1}{3}\int_0^\infty \langle\vec{v}(0)\cdot\vec{v}(t)\rangle\,dt.
\]

\section*{Exemplos Resolvidos}

\subsection*{Exemplo 1 --- Coeficiente de difus\~ao da \'agua}
DM da \'agua l\'iquida a 300 K. MSD medido: ap\'os $t=5$ ps, MSD $= 0{,}15$ nm$^2$.
$D = 0{,}15\times10^{-18}/(6\times5\times10^{-12}) = 5\times10^{-9}$ m$^2$/s.
Valor experimental de $D$ para \'agua: $\approx 2{,}3\times10^{-9}$ m$^2$/s. (Diferen\c{c}a indica limitac\~ao do campo de for\c{c}a.)

\subsection*{Exemplo 2 --- Interpreta\c{c}\~ao da RDF}
Para LJ a $T^*=0{,}7$, $\rho^*=0{,}8$ (l\'iquido): $g(r)$ tem pico principal em $r\approx1{,}1\sigma$ (vizinhos mais pr\'oximos), pico secund\'ario em $r\approx2\sigma$ (segunda camada), e $g\to1$ para $r>4\sigma$.

\section*{Exerc\'icios}

\begin{enumerate}
  \item \textbf{(MSD)} Uma part\'icula LJ percorre as dist\^ancias: $d(0)=0$, $d(1)=0{,}15$, $d(2)=0{,}28$, $d(3)=0{,}45$ (em $\sigma^2$, unidades reduzidas de tempo). Estime $D$ ajustando MSD $= 6Dt$.
  
  \item \textbf{(RDF)} Descreva o padr\~ao esperado de $g(r)$ para: (a) g\'as de LJ; (b) l\'iquido de LJ; (c) s\'olido cristalino de LJ. Quais as caracter\'isticas distintas de cada fase?
  
  \item \textbf{(N\'umero de coordena\c{c}\~ao)} O n\'umero de coordena\c{c}\~ao \'e $n_c = 4\pi\rho\int_0^{r_\text{min2}} g(r)r^2\,dr$. Para argônio l\'iquido ($\rho^*=0{,}85$), $\int_0^{r_\text{min2}} g(r)r^2\,dr \approx 2{,}9$. Calcule $n_c$.
  
  \item \textbf{(Regime bal\'istico)} Para $t < \tau$ (tempo livre m\'edio), MSD $\approx \langle v^2\rangle t^2 = 3k_BT/m \cdot t^2$. Para LJ a $T^*=1{,}5$ (em unidades reduzidas), estime o MSD em $t^*=0{,}1$.
  
  \item \textbf{(Subdifus\~ao)} Uma prote\'ina numa membrana biol\'ogica mostra MSD $\sim t^{0{,}7}$. O que isso indica? Como isso difere de difus\~ao normal?
  
  \item \textbf{(VACF)} A fun\c{c}\~ao de autocorrela\c{c}\~ao de velocidades $C(t) = \langle\vec{v}(0)\cdot\vec{v}(t)\rangle$ decai exponencialmente em gases ($C(t) \sim e^{-t/\tau}$) e oscila em s\'olidos. O que $C(0)$ representa? Como $C(t)$ se comporta para um s\'olido cristalino?
  
  \item \textbf{(Transi\c{c}\~ao de fase)} Ao aumentar a temperatura em $V$ fixo, como $g(r)$ e $D$ mudam durante a transi\c{c}\~ao s\'olido-l\'iquido? O que \'e uma descontinuidade em $D$ vs.\ $T$?
\end{enumerate}

\section*{Resumo R\'apido}
\begin{itemize}
  \item RDF: $g(r)$ mostra estrutura local; picos = camadas de vizinhos.
  \item MSD $= 6Dt$ (3D, difus\~ao normal); $D$ = inclinac\~ao/6.
  \item Green-Kubo: $D = \frac{1}{3}\int\langle\vec{v}(0)\cdot\vec{v}(t)\rangle dt$.
\end{itemize}

\section*{Refer\^encias}
\begin{itemize}
  \item Frenkel, D.; Smit, B.\ \textit{Understanding Molecular Simulation}. Academic Press, 2002.
  \item Allen, M.\ P.; Tildesley, D.\ J.\ \textit{Computer Simulation of Liquids}. Oxford, 2017.
  \item Notas de aula da disciplina.
\end{itemize}

\end{document}
